% ============================================================
% Study 2 — Calibration across prediction horizons
% ============================================================

\subsection*{Study 2}

Study~1 established overconfidence at a single, short horizon. Study~2 asked
whether stated confidence intervals---which must widen to maintain coverage
as forecast uncertainty grows---actually scale with prediction horizon.
Using a backtesting design with 100 large-cap equities and nine horizons spanning
1 to 22 trading days, we tested whether miscalibration grows, shrinks, or
diverges across models as the horizon extends.

\subsubsection*{Methods}

The same three models were evaluated using the same prompt format and
structured JSON output as Study~1. Study~2 used a backtesting design in which
predictions were made against historical prices using date-blinded prompts
(current price and horizon provided; calendar date withheld) to prevent
models from conditioning on known future events.

\paragraph{Design} One hundred large-cap S\&P~500 equities were selected.
Predictions were collected at nine horizons: 1, 2, 3, 6, 7, 18, 20, 21, and
22 calendar days. Five independent predictions per ticker per model were
collected at each horizon, using a reference date of March~24, 2026.
Market closures were handled by applying a $\pm$6-calendar-day tolerance
window when matching predictions to realized prices. After excluding
unresolved predictions, the design yielded between 157 and 1,495 scored
predictions per model-horizon cell ($N > 18{,}000$ total).

\paragraph{Calibration metrics} Hit rates and ECE were computed as in
Study~1, aggregated within each model-horizon cell.

\subsubsection*{Results}

Table~\ref{tab:study2} reports 80\% hit rates and ECE at four representative
horizons; complete results for all nine horizons are in the Web Appendix.

\paragraph{Horizon degradation: Claude and GPT-4} Both models exhibited
marked degradation with increasing horizon. Claude's Hit@80 fell from
80.5\% at 1~day to 46.2\% at 18~days (ECE rising from 1.4\% to 29.8\%),
consistent with CI widths anchored to short-horizon volatility that fail
to scale with realized uncertainty at longer windows. GPT-4---already
poorly calibrated at 1~day (ECE\,=\,21.8\%)---deteriorated further, reaching
Hit@80\,=\,27.2\% and ECE\,=\,47.4\% at 18~days. By 22~days, GPT-4's ECE
(49.9\%) approached the theoretical maximum for three-level calibration.

\paragraph{Horizon robustness: Gemini} Gemini exhibited a qualitatively
different pattern. Its 1-day calibration (ECE\,=\,1.8\%, Hit@80\,=\,78.5\%) was
the best of any model at any single horizon in the study. While it degraded
at intermediate horizons (ECE\,=\,16.6\% at 18~days), it recovered strongly
at 22~days (Hit@80\,=\,80.1\%, ECE\,=\,2.5\%)---near-perfect calibration at the
longest horizon tested.

\begin{table}[ht]
  \centering
  \caption{Study~2: Hit@80 and ECE by model and prediction horizon
    (selected horizons; full table in Web Appendix).}
  \label{tab:study2}
  \begin{tabular}{lcccccc}
    \toprule
    & \multicolumn{2}{c}{Claude}
    & \multicolumn{2}{c}{Gemini}
    & \multicolumn{2}{c}{GPT-4} \\
    \cmidrule(lr){2-3}\cmidrule(lr){4-5}\cmidrule(lr){6-7}
    Horizon & H@80   & ECE    & H@80   & ECE    & H@80   & ECE    \\
    \midrule
    1d  & 80.5\% & 1.4\%  & 78.5\% & 1.8\%  & 55.2\% & 21.8\% \\
    6d  & 56.0\% & 20.6\% & 71.7\% & 8.4\%  & 34.5\% & 41.6\% \\
    18d & 46.2\% & 29.8\% & 59.8\% & 16.6\% & 27.2\% & 47.4\% \\
    22d & 51.2\% & 26.9\% & 80.1\% & 2.5\%  & 24.4\% & 49.9\% \\
    \bottomrule
  \end{tabular}
\end{table}

\subsubsection*{Discussion}

Prediction horizon strongly moderated miscalibration in a model-specific
pattern. The degradation of Claude and GPT-4 mirrors the human tendency
toward greater interval underestimation at longer horizons \citep{Griffin1992},
consistent with CI widths anchored to short-run volatility. Gemini's
recovery at 22~days is more surprising: its CI widths appear to scale
approximately proportionally to realized volatility at the extremes of the
tested range but not at intermediate horizons. Whether this reflects a genuine
uncertainty representation or an incidental match with equity return
volatility at these specific windows requires replication at additional
horizons and across different historical periods.
